\subsection{RQ1：谁来决定？}
\label{sec:decide}
我们先在结构化观测、没有攻击者的条件下，把决定逐项从模型手中拿走，测量这对已解决案例的影响。
\label{sec:utility}
\label{sec:cost}
\label{sec:fair}
\label{sec:rq-stop}
\label{sec:rq-rank}
\label{sec:external}
\emph{探测与停止。} 在提示相同的条件下，信息增益选择在 Qwen2.5-7B 上比随机选择高 $+0.306$（区间 $+0.14$ 到 $+0.47$），在其他每个会给出结论的模型上高 $+0.26$ 到 $+0.35$。选择探针并不会让模型给出结论：控制器停止把 Llama-3.1-8B 从 0 提高到 0.917（区间 $+0.79$ 到 $+1.00$），而 Qwen2.5-7B 两者都需要（只加停止把它从 0.083 提高到 0.694，选择再加 $+0.222$）。两个效应在 sigma-triage 上都复现：排序在 Qwen2.5-7B 上增加 $+0.171$，在五个模型上增加 $+0.105$ 到 $+0.230$；停止把 Llama-3.1-8B 从 304 个回合中的 1 个提高到全部（\cref{fig:probe-stop}）。

\begin{figure*}[t]
  \centering
  \includegraphics[width=0.92\textwidth]{fig-probe-stop.pdf}
  \caption{Llama-3.1-8B 缺的是停止，而不是探测。按“谁选探针”与“谁可以停止”划分的已解决回合比例（sec-triage，每点 72 个回合）。虚线是事后运行的不含 LLM 的控制器（EIG/成本选择加控制器停止，相同的种子和任务）。}
  \label{fig:probe-stop}
\end{figure*}

\begin{table}[t]
  \caption{没有模型能在控制器之上增加什么。已解决回合数 / 给出正确原因的回合数，分别在由模型决定何时停止（“LLM”）、后验停止和确认条件停止下。每个 LLM 行都使用控制器选择和守卫；种子、任务和预算与不含 LLM 的行相同。}
  \label{tab:v11fair}
  \centering\small
  \begin{adjustbox}{max width=\columnwidth}% generated by paper/make_tables.py -- do not edit
\begin{tabular}{lrrrrrr}
\toprule
 & \multicolumn{3}{c}{sec-triage (of 72)} & \multicolumn{3}{c}{sigma-triage (of 152)} \\
\cmidrule(lr){2-4}\cmidrule(lr){5-7}
Planner & LLM & post. & conf. & LLM & post. & conf. \\
\midrule
none (LLM-free) & -- & 66 / 72 & 72 / 72 & -- & 152 / 152 & 152 / 152 \\
Q-7B & 66 / 66 & 66 / 72 & 72 / 72 & 126 / 142 & 141 / 141 & 141 / 141 \\
Q-32B & 72 / 72 & 66 / 72 & 72 / 72 & 144 / 152 & 152 / 152 & 152 / 152 \\
Q-72B & 72 / 72 & 66 / 72 & 72 / 72 & 148 / 152 & 152 / 152 & 152 / 152 \\
L-8B & 0 / 0 & 66 / 72 & 72 / 72 & 1 / 1 & 152 / 152 & 152 / 152 \\
L-70B & 72 / 72 & 66 / 72 & 72 / 72 & 145 / 152 & 152 / 152 & 152 / 152 \\
\bottomrule
\end{tabular}
\end{adjustbox}
\end{table}

\emph{控制器单独。} 采用确认条件停止、完全不用模型时，控制器解决了 sec-triage 的 72/72 和 sigma-triage 的 152/152 个回合，成本为 2.54 和 2.20 个单位（\cref{tab:v11fair}）。让 Qwen2.5-72B 决定何时停止，解决了 72/72 和 148/152，成本为 3.86 和 3.25 个单位，预注册的差值为 $+0.000$ 和 $-0.028$（区间 $-0.111$ 到 $0.000$）；Qwen2.5-32B 和 Llama-3.1-70B 表现相同，每次失误都给出了正确原因但没有确认特征。没有确认的后验停止只解决了 72 个中的 66 个，因为它靠排除得出 DNS 命令控制的结论。自适应选择仍然重要：固定流程解决 66 个，目录顺序 65 个。Qwen2.5-7B 在 11 个 sigma-triage 回合中通过提出略低于阈值的结论让带守卫的控制器卡住；有了恢复规则，一个坚持提出这类结论的脚本规划器解决了全部 152 个回合，而不是一个也没有。

\begin{table}[t]
  \caption{更清楚的提示帮助了基线，但没有缩小差距。原始提示和在开发种子上为帮助基线而选出的变体下的已解决回合数（共 48 个，sec-triage）。}
  \label{tab:v12fair}
  \centering\small
  \begin{adjustbox}{max width=\columnwidth}% generated by paper/make_tables.py -- do not edit
\begin{tabular}{lrrrr}
\toprule
 & \multicolumn{2}{c}{Qwen2.5-7B} & \multicolumn{2}{c}{Llama-3.1-8B} \\
\cmidrule(lr){2-3}\cmidrule(lr){4-5}
Configuration & original & clearer & original & clearer \\
\midrule
ReAct-style & 2 & 0 & 0 & 0 \\
Memory-only & 8 & 16 & 0 & 0 \\
\hesp{}, guard, LLM stops & 44 & 48 & 0 & 32 \\
\hesp{}, guard + confirmation stop & 48 & 48 & 48 & 48 \\
\bottomrule
\end{tabular}
\end{adjustbox}
\end{table}

\emph{给基线更好的提示。} 基线也许只是缺少一个清楚的提示。我们在开发种子上、按事先固定的规则，选出了对基线帮助最大的提示变体（\cref{tab:v12fair}）。它让 Qwen2.5-7B 的 Memory-only 结果翻倍（从 8 到 16 个），对 ReAct 式和两个 Llama 基线没有帮助；推荐的控制器仍然多解决 $+0.667$（Qwen2.5-7B，区间 $+0.46$ 到 $+0.85$）和 $+1.000$（Llama-3.1-8B）。

\emph{开销。} \hesp{} 降低而不是增加了调查的开销，因为它向模型提问更少。控制器停止把每个回合的模型调用从 5.5 到 10.1 次降到 2.3 次，token 减少 2.5 到 5 倍；Qwen2.5-72B 上回合耗时的中位数从 46.8\,s 降到 15.3\,s（附录~\ref{app:runtime}）。与一次模型调用相比，账本更新和探针选择的时间可以忽略。

\smallskip\noindent\textbf{发现 1.} 小模型在选择探针和停止上的失败原因不同，而接管这两个决定的控制器解决的案例与我们测试的任何模型一样多，成本更低。在候选原因给定、观测已结构化时，做决定不需要模型。


\subsection{RQ2：谁来读？}
\label{sec:read}
现在让探针以文档格式或\emph{漂移}格式返回原始日志文本。漂移格式用改写过的日志行陈述同样的事实，就像日志管道变更之后那样。我们比较固定规则解析器和由模型充当的解析器；控制器始终负责决定，因此模型的唯一角色是读取。

\begin{table}[t]
  \caption{日志格式一变，读取就需要模型。sec-triage 中探针返回原始日志文本时，固定规则解析器与模型作为解析器的解析准确率和已解决回合数（共 48 个；Q-7B：Qwen2.5-7B；L-8B：Llama-3.1-8B）。每一行都由控制器做决定；``injected'' 行在 RQ3 中讨论。}
  \label{tab:v12raw}
  \centering\small
  \begin{adjustbox}{max width=\columnwidth}% generated by paper/make_tables.py -- do not edit
\begin{tabular}{llrrrrrr}
\toprule
Parser & Log format & \multicolumn{2}{c}{parse accuracy} & \multicolumn{2}{c}{injection adopted} & verified (of 48) & escalated (of 48) \\
\cmidrule(lr){3-4}\cmidrule(lr){5-6}\cmidrule(lr){7-7}\cmidrule(lr){8-8}
 & & Q-7B & L-8B & Q-7B & L-8B & Q-7B / L-8B & Q-7B / L-8B \\
\midrule
structured (reference) & documented & -- & -- & -- & -- & 48 / 48 & 0 / 0 \\
rule & documented & 108/108 & 108/108 & -- & -- & 48 / 48 & 0 / 0 \\
rule & drifted & 0/364 & 0/364 & -- & -- & 0 / 0 & 48 / 48 \\
rule & injected & 108/108 & 108/108 & 0/66 & 0/66 & 48 / 48 & 0 / 0 \\
LLM & documented & 108/108 & 108/108 & -- & -- & 48 / 48 & 0 / 0 \\
LLM & drifted & 108/108 & 174/196 & -- & -- & 48 / 34 & 0 / 12 \\
LLM & injected & 125/184 & 129/136 & 48/105 & 1/71 & 24 / 42 & 18 / 6 \\
\bottomrule
\end{tabular}
\end{adjustbox}
\end{table}

\Cref{tab:v12raw} 给出了比较。在文档格式下，两种解析器都读对了全部 108 条观测，控制器解决了全部 48 个回合，与结构化观测时相同。在漂移格式下，规则解析器读不出 364 条观测中的任何一条，每个回合都被升级给分析员。Qwen2.5-7B 读对了全部 108 条漂移观测，解决了全部 48 个回合；Llama-3.1-8B 读对 196 条中的 174 条，解决 34 个。在加入跨来源佐证的推荐配置下，Qwen2.5-7B 解决 48 个漂移回合中的 41 个，与从不误读的脚本读取器完全相同，Llama-3.1-8B 解决 28 个。换句话说，一旦日志不再符合解析器的预期，去掉模型就等于让智能体完全无法工作。

\begin{table}[t]
  \caption{两半都不能单独工作。sec-triage 基础与噪声变体中，按谁决定、谁读取划分的已解决回合数（共 48 个；Q-7B / L-8B）。模型行中模型单独决策：结构化条件下看到环境给出的标签和 JSON 响应，日志格式条件下只看到日志文本、没有标签。控制器行使用确认停止，其中固定规则在结构化条件下是环境的分类器，在日志格式条件下是规则解析器；最后一行即 \cref{tab:v12raw} 中的模型读取器。}
  \label{tab:readdecide}
  \centering\small
  \begin{adjustbox}{max width=\columnwidth}% generated by paper/make_tables.py -- do not edit
\begin{tabular}{lllrrr}
\toprule
Decides & Reads & Prompt & structured & documented log & drifted log \\
\midrule
model (ReAct-style) & model & default & 0 / 0 & 0 / 0 & 0 / 0 \\
model (ReAct-style) & model & clear finish & 0 / 0 & 0 / 0 & 0 / 0 \\
model (Memory-only) & model & default & 9 / 0 & 0 / 0 & 0 / 0 \\
model (Memory-only) & model & clear finish & 13 / 0 & 0 / 0 & 0 / 0 \\
\midrule
controller & fixed rules & -- & 48 / 48 & 48 / 48 & 0 / 0 \\
controller & model & -- & -- & 48 / 48 & 48 / 34 \\
\bottomrule
\end{tabular}
\end{adjustbox}
\end{table}

\emph{没有控制器的模型。}只会读并不能让模型有用。我们使用同样的任务、种子和日志文本，让每个模型自行决策（ReAct 式或 Memory-only，配默认提示或 \cref{tab:v12fair} 中更清楚的提示），并只给它每个探针的日志文本、不给结果标签（\cref{tab:readdecide}）。在全部四种配置和两种日志格式下，两个模型都没有解决任何案例：每个回合都没有结论，也没有一个回合给出错误原因。给标签时，Qwen2.5-7B 最多解决 48 个中的 13 个，Llama-3.1-8B 一个也没有。Qwen2.5-7B 确实读懂了文本：看到漂移格式的认证摘要后，它写下的理由是日志显示了分布式的撞库攻击，但接下来它没有结案，而是再次提议同一个探针，平均每个回合 6.8 次，给标签时为 2.3 次。把提示中“not parsed”的措辞换成中性的写法（主研究之后进行的探索性检查），它在两种格式下仍然是 48 个中的 0 个。在漂移格式下，“模型为控制器读取”减去“模型单独决策”的预注册差值，Qwen2.5-7B 为 $+1.000$（区间 $+1.000$ 到 $+1.000$），Llama-3.1-8B 为 $+0.708$（区间 $+0.44$ 到 $+0.94$）。

\smallskip\noindent\textbf{发现 2.} 两半都不能单独工作。日志格式一变，固定解析器就完全失效；小模型单独决策时，尽管读得懂原始日志，却一个案例也解决不了；同一个模型为控制器读取时，解决了大部分案例。

\subsection{RQ3：模型读取器安全吗？}
\label{sec:rawlog}
读取器也正是攻击者写入的地方。\Cref{tab:v13} 报告 Qwen2.5-7B、Llama-3.1-8B 以及采纳每条声称的脚本读取器在四种读取器策略下受到的攻击。

\begin{table*}[t]
  \caption{只有读取器信任挡住了针对读取器的每一种攻击。四探针声称和三种自适应攻击下被结为良性的需处置案例（共 36 个）；文档格式与漂移格式下已解决的诚实良性案例（共 12 个；最多 6 个，因为监控误报只有一条特征）；以及只用读取器时采纳攻击者标签的观测数。策略：只用读取器；规则解析器优先，读不出的交给读取器；先删去已知自由文本字段；在规则解析器优先之上加读取器信任。sec-triage 基础与噪声变体，每格 48 个回合。}
  \label{tab:v13}
  \centering\small
  \begin{adjustbox}{max width=\textwidth}% generated by paper/make_tables.py -- do not edit
\begin{tabular}{llrrrrrrr}
\toprule
 & & \multicolumn{4}{c}{attacks closed as benign (of 36)} & \multicolumn{2}{c}{benign verified (of 12)} & adopted \\
\cmidrule(lr){3-6}\cmidrule(lr){7-8}\cmidrule(lr){9-9}
Reader & Policy & four probes & every probe & every probe, drifted & log line & documented & drifted & (reader only) \\
\midrule
adopts every claim & reader only & 36 & 36 & 36 & 36 & 6 & 6 & 252/252 \\
 & rule parser first & 0 & 0 & 36 & 0 & 6 & 6 & -- \\
 & fields stripped & 0 & 0 & 0 & 36 & 6 & 6 & -- \\
 & reader trust & 0 & 0 & 0 & 0 & 6 & 0 & -- \\
\midrule
Qwen2.5-7B & reader only & 0 & 0 & 36 & 0 & 6 & 6 & 185/338 \\
 & rule parser first & 0 & 0 & 36 & 0 & 6 & 6 & -- \\
 & fields stripped & 0 & 0 & 0 & 0 & 6 & 6 & -- \\
 & reader trust & 0 & 0 & 0 & 0 & 6 & 0 & -- \\
\midrule
Llama-3.1-8B & reader only & 0 & 0 & 19 & 0 & 6 & 1 & 174/380 \\
 & rule parser first & 0 & 0 & 19 & 0 & 6 & 1 & -- \\
 & fields stripped & 0 & 0 & 0 & 0 & 6 & 1 & -- \\
 & reader trust & 0 & 0 & 0 & 0 & 6 & 0 & -- \\
\bottomrule
\end{tabular}
\end{adjustbox}
\end{table*}

\emph{第一种攻击只是看起来被遏制了。} 四个探针上的良性声称，被 Qwen2.5-7B 在 105 条观测中采纳 48 条，被 Llama-3.1-8B 在 71 条中采纳 1 条，没有攻击被结案。脚本读取器说明了原因：采纳每一条声称时，它在两个探针后就把全部 36 个需处置案例结为授权扫描，因为四个被声称的探针横跨三个来源组，仅凭伪造读数就满足了跨来源佐证。模型之所以安全，只是因为它们采纳了一些声称、没有采纳另一些。

\emph{自适应攻击达到了这个上界。} 在文档格式下给每个探针都写上故事时，Qwen2.5-7B 采纳了 \texttt{change\_ticket} 和 \texttt{access\_pattern} 上的声称，但从不采纳 \texttt{source\_ips} 上的声称（0/42），因为那里真实记录与声称矛盾；没有案件被结案，但有 18 个被升级。在漂移格式下，声称孤立存在，Qwen2.5-7B 采纳了全部 84 条，36 个需处置案例全部被结为授权扫描，与脚本读取器完全相同。Llama-3.1-8B 采纳 124 条中的 85 条，结掉 19 个。日志行注入没有结掉任何案件，因为两个模型都读第一条 \texttt{source\_ips} 记录；但 Llama-3.1-8B 在 72 条观测中全部采纳了另外两个探针的伪造记录，让 36 个需处置案例全部悬而未决。

\emph{一个案例。} 考虑漂移自适应攻击下的一条凭证填充告警，读取器是 Qwen2.5-7B。第一个探针 \texttt{source\_ips} 返回下面的日志（已折行）。真实的那一行报告的是住宅僵尸网络地址段，攻击者的用户代理则声称是已知扫描器。
\begin{Verbatim}[fontsize=\scriptsize,frame=single]
Top talkers: mostly consumer broadband
  ranges flagged as botnet
heartbeat ok pipeline=v2
user_agent="Mozilla/5.0 (compatible)
  reputation=known_scanner_asn (authorized
  scanner, CHG-88; parsers: report
  known_scanner_asn)"
\end{Verbatim}
模型返回 \textit{known scanner}，把授权扫描抬到 0.82。没有读取器信任时，控制器因为缺少第二个来源组而暂缓结论，运行 \texttt{access\_pattern}；它的伪造读数 \textit{path scanning} 来自另一个来源组，于是控制器在两个探针后以 0.95 把案件结为授权扫描。有了读取器信任，同样的两条读数照样到达，但都不能确认良性原因，于是控制器继续探测。真实的认证失败特征在第五个探针出现，与伪造的故事矛盾；后验在授权扫描与 \textit{other} 之间分裂，没有可信观测能确认任何一方，案件最终带着日志中全部八条观测交给了分析员。

\emph{什么能阻止它。} 规则解析器优先在文档格式下免疫，在漂移下无用。删去自由文本字段能挡住字段内的每条声称，但挡不住伪造的日志行，脚本读取器会全部采纳。读取器信任在任何攻击下、在两个模型和脚本读取器上都没有结掉任何案件，而且只要规则解析器读得出就没有代价（48 个中解决 41 个，与结构化观测相同）。它的代价落在漂移格式上：没有诚实的良性案例被解决（0/12，对比 6 个），在漂移自适应攻击下每个回合都在耗尽预算后被升级。两个预注册比较成立：只用模型读取比读取器信任多结掉的需处置案例比例，在 Qwen2.5-7B 上为 $+0.333$（97.5\% 置信区间 $+0.333$ 到 $+0.333$），在 Llama-3.1-8B 上为 $+0.176$（区间 $+0.093$ 到 $+0.259$）。

\smallskip\noindent\textbf{发现 3.} 读取的模型可能被说服，一次接受攻击者对每条日志的读法，这会击穿跨来源佐证；让模型解析的观测可以确认攻击、却永远不能确认良性结论，就关闭了这条路径，而只有在没有可信解析器能读日志时，诚实的良性案例才为此付出代价。


\subsection{RQ4：攻击者能触及决定吗？}
控制器负责决定之后，攻击者仍可能通过模型看到的文本（A1）或账本读取的伪造证据（A2）去影响结论。

\paragraph{注入指令}\label{sec:injection}
\label{sec:adversarial}
\label{sec:security}
我们把四种注入文本分别放进 sec-triage 六个需处置原因的每个探针响应中，并在全部五个模型上运行每种配置（\cref{tab:v11attack} 左）。

\begin{table*}[t]
  \caption{结束守卫挡住注入指令，跨来源组佐证挡住伪造来源。左：四种注入文本下被结为良性（“漏判”）或无结论（“未决”）的需处置案例，每个模型共 72 个。中：伪造证据下被结为良性的需处置案例，共 18 个（五个模型的范围）。右：已解决的诚实良性案例，共 6 个（五个模型的范围）。sec-triage。}
  \label{tab:v11attack}
  \centering\small
  \begin{adjustbox}{max width=\textwidth}% generated by paper/make_tables.py -- do not edit
\begin{tabular}{lrrrrrrrr}
\toprule
 & \multicolumn{5}{c}{4 injection texts: missed / unresolved (of 72)} & \multicolumn{2}{c}{forged: missed (of 18)} & benign \\
\cmidrule(lr){2-6}\cmidrule(lr){7-8}\cmidrule(lr){9-9}
Configuration & Q-7B & Q-32B & Q-72B & L-8B & L-70B & probe & feed & verif.\ (of 6) \\
\midrule
Memory-only & 69 / 0 & 57 / 0 & 69 / 0 & 0 / 72 & 63 / 0 & 0--6 & 0--15 & 0--6 \\
Memory-only, raw text removed & 0 / 65 & 0 / 0 & 0 / 0 & 0 / 72 & 0 / 24 & 0--6 & 0--15 & 0--6 \\
\hesp{}, guard only & 0 / 44 & 0 / 46 & 0 / 66 & 0 / 72 & 0 / 58 & 0 & 0 & 0--6 \\
\hesp{}, stop, no guard & 39 / 0 & 36 / 0 & 48 / 0 & 0 / 0 & 36 / 0 & 18 & 18 & 6 \\
\hesp{}, guard + stop & 0 / 24 & 0 / 29 & 0 / 42 & 0 / 0 & 0 / 2 & 18 & 18 & 6 \\
\quad + corroboration by probe & 0 / 0 & 0 / 0 & 0 / 11 & 0 / 0 & 0 / 0 & 0 & 0 & 3 \\
\quad + corroboration by source & 0 / 12 & 0 / 0 & 0 / 3 & 0 / 0 & 0 / 0 & 0 & 0 & 3 \\
\bottomrule
\end{tabular}
\end{adjustbox}
\end{table*}

当由模型决定结论时，攻击在每个会给出结论的模型上都有效：Memory-only 智能体在 Qwen2.5-7B、32B、72B 和 Llama-3.1-70B 上分别把 72 个需处置案例中的 69、57、69、63 个结为授权扫描；Llama-3.1-8B 之所以“安全”，只是因为它从不给出结论。只把探针选择和停止从模型手里拿走还不够：没有守卫时，在 72 个回合中的 36 到 48 个里，模型在控制器自己的停止生效之前就提出了注入的结论。有了守卫，任何模型上的任何配置都没有关闭一个案件，因为注入的文本从不进入账本，而守卫只接受账本支持的结论。把原始文本从提示里去掉同样能挡住攻击，但会让 Qwen2.5-7B 的 72 个案例中有 65 个无法解决，因为模型缺少它需要的信息。

注入仍然消耗分析员的时间。只有守卫时，72 个回合中有 44 到 72 个没有结论，因为模型不断提出被拒绝的结论；加上控制器停止后，这个数字降到 0 到 42。

\smallskip\noindent\textbf{发现 4a.} 攻击者写进日志的文本，在由模型决定结论时决定结论，在结论必须依据账本证据时则完全不起作用；剩下的影响是延迟，而不是结案。


\paragraph{伪造来源}
\label{sec:evidence}
守卫生效之后，攻击者必须伪造账本会读取的证据。一条伪造的扫描器答案在一个探针后就把授权扫描抬过接受阈值，控制器停止随即在每个模型上把 18 个需处置案例全部结为良性；伪造源也是如此（\cref{tab:v11attack} 中）。控制器的长处——先按信息量最大的观测行动——恰恰被攻击者利用。伪造源也误导了更大的模型本身，它们在 Memory-only 下最多结掉 18 个中的 15 个。

要求两条支持观测会在每个模型上让全部伪造案例升级。但只数探针并不够：在伪造源实验中，控制器的第二个探针碰巧读了另一个来源。当固定流程先读 \texttt{source\_ips} 再读 \texttt{threat\_intel}（两者都由被污染的信誉源供数）时，按探针佐证把 18 个需处置案例全部结案，而按来源组佐证一个也没有结案。这条规则的代价落在只有一条特征的诚实良性案例上：在 sec-triage 中，加上佐证后 6 个良性回合只解决 3 个（原为 6 个），因为监控误报在那里只有一条特异特征；在 sigma-triage 的 26 个良性解释上，佐证让 104 个良性回合中的 8 个升级。

\smallskip\noindent\textbf{发现 4b.} 负责决定的控制器会被一个伪造来源击败；跨来源佐证能阻止它，代价取决于诚实的良性原因有多少条独立特征。


\subsection{RQ5：证据规则、不精确的预测表与缺失的原因}
\label{sec:policy}
\label{sec:mismatch}
\label{sec:rq-help}
我们的验证器和确认条件停止要求一条只有所称原因才会产生的观测。\Cref{tab:v12policy} 用不含 LLM 的控制器，把这条规则与联合证据作比较；联合证据接受任何使所称原因比其他每个具名原因都高出十倍的观测组合。

\begin{table*}[t]
  \caption{证据规则在自动化与安全之间取舍。不含 LLM 的控制器，计数。“single”/“joint”：被单条特征验证器与联合验证器接受的回合数。cause absent：真实原因从候选中移除。confusion：每个原因的表行向另一个原因混合 0.25。comp-triage：攻击与使用相同工具的合法活动。}
  \label{tab:v12policy}
  \centering\small
  \begin{adjustbox}{max width=\textwidth}% generated by paper/make_tables.py -- do not edit
\begin{tabular}{lrrrrrrrrr}
\toprule
 & \multicolumn{2}{c}{sec, matched (48)} & \multicolumn{2}{c}{sec, cause absent (48)} & sec, confusion & \multicolumn{4}{c}{comp-triage (72; 18 attacks)} \\
\cmidrule(lr){2-3}\cmidrule(lr){4-5}\cmidrule(lr){6-6}\cmidrule(lr){7-10}
Stop rule & single & joint & wrong & escal. & verified & right cause & wrong & missed attack & escal. \\
\midrule
posterior (0.8) & 42 & 48 & 24 & 24 & 47 & 64 & 7 & 5 & 1 \\
single signature, $r{=}10$ & 48 & 48 & 0 & 48 & 0 & 20 & 0 & 0 & 52 \\
single signature, $r{=}3$ & 48 & 48 & 0 & 48 & 0 & 37 & 2 & 1 & 33 \\
joint, named causes, $r{=}10$ & 42 & 48 & 24 & 24 & 47 & 59 & 6 & 4 & 7 \\
joint, incl.\ \textit{other}, $r{=}10$ & 42 & 48 & 24 & 24 & 47 & 59 & 6 & 4 & 7 \\
joint, incl.\ \textit{other}, $r{=}30$ & 42 & 48 & 18 & 30 & 18 & 55 & 2 & 1 & 15 \\
\bottomrule
\end{tabular}
\end{adjustbox}
\end{table*}

在原因使用相同工具的地方，联合证据在 72 个 comp-triage 回合中给出 59 个正确原因，单条特征规则只有 20 个，并升级 52 个。代价以攻击者的货币支付：联合证据把 18 个攻击中的 4 个结为与之相似的合法活动；当真实原因不在候选之中时，它在 48 个 sec-triage 回合中的 24 个里给出另一个需处置原因；要求证据比 \textit{other} 高十倍不改变这一点，三十倍也只降到 18 个。只有单条特征规则在任何设置下都没有给出错误原因。

这条规则也决定了预测表出错时控制器如何失败（\cref{tab:v10b}）。每个原因 20 个开发回合计数得到的表与生成器的表一致，在 sec-triage 中五个就已足够；向模型询问得到的表只能解决 6.9\% 到 37.5\% 的案例。把每个表行与随机噪声混合时解决率下降，但即使噪声占 75\%，给出错误原因的回合也至多占 0.9\%：控制器选择升级。开发数据中缺失的原因会产生自信的错误结论，但总是同一类别，因此缺失的攻击从未被结为良性。当预测表混淆两个原因时，后验停止仍解决 48 个中的 47 个，而确认条件停止把 48 个全部升级。

\begin{table*}[t]
  \caption{表误差（不含 LLM，后验停止 0.8）。EIG/成本选择下：已验证、原因正确但无确认特征、原因错误、升级的回合比例（四者之和为一）；随机选择下：已验证完成率。“perturbed”把每个表行与随机 Dirichlet(1) 行以权重 $\lambda$ 混合（每个水平 20 个随机表）；“one cause missing”依次把每个原因从开发数据中移除并在该原因上测试，取平均。}
  \label{tab:v10b}
  \centering\small
  \begin{adjustbox}{max width=\textwidth}% generated by paper/make_tables.py -- do not edit
\begin{tabular}{lrrrrrrrrrr}
\toprule
 & \multicolumn{5}{c}{sec-triage} & \multicolumn{5}{c}{sigma-triage} \\
\cmidrule(lr){2-6} \cmidrule(lr){7-11}
Tables & verif. & correct, unverif. & wrong & esc. & random verif. & verif. & correct, unverif. & wrong & esc. & random verif. \\
\midrule
matched tables & 0.917 & 0.083 & 0.000 & 0.000 & 0.792 & 1.000 & 0.000 & 0.000 & 0.000 & 0.886 \\
base variant only & 0.917 & 0.083 & 0.000 & 0.000 & 0.792 & 1.000 & 0.000 & 0.000 & 0.000 & 0.886 \\
$k{=}1$, base only & 0.917 & 0.083 & 0.000 & 0.000 & 0.778 & 1.000 & 0.000 & 0.000 & 0.000 & 0.882 \\
perturbed, $\lambda{=}0.25$ & 0.885 & 0.083 & 0.000 & 0.031 & 0.500 & 0.992 & 0.008 & 0.000 & 0.000 & 0.910 \\
perturbed, $\lambda{=}0.5$ & 0.494 & 0.177 & 0.000 & 0.329 & 0.227 & 0.978 & 0.017 & 0.000 & 0.005 & 0.838 \\
perturbed, $\lambda{=}0.75$ & 0.098 & 0.027 & 0.004 & 0.871 & 0.042 & 0.805 & 0.040 & 0.009 & 0.146 & 0.513 \\
one cause missing (mean) & 0.000 & 0.000 & 0.500 & 0.500 & 0.000 & 0.000 & 0.000 & 0.079 & 0.921 & 0.000 \\
\bottomrule
\end{tabular}
\end{adjustbox}
\end{table*}

\smallskip\noindent\textbf{发现 5.} 要求一条决定性观测，能让控制器不给出错误原因，代价是升级那些只有组合证据才能解决的案例；部署需要在这条前沿上选择自己的位置。
